% Chunk: Chapter 7 Appendix---Dirac Brackets from Canonical Commutators
% Source: Weinberg QFT Vol. I, physical PDF pp. 355--360 (printed pp. 332--337)
% Figures: none.
% Uncertainties: none.

\chapterappendix{}{Dirac Brackets from Canonical Commutators}

In this Appendix we shall show, in theories of two types, that the formula
giving commutators as Dirac brackets times $i$ follows from the usual
canonical commutation relations for a reduced set of variables.

\begin{center}
$A$
\end{center}

Suppose (as in the case of a massive vector field $V^\mu$) that the quantum
variables $\Psi^a$ and $\Pi_a$ appearing in the Lagrangian $L$ may be divided
into two classes:\footnote{We are again using a compact notation, in which
labels like $a$, $n$, and $r$ include a space coordinate $\mathbf{x}$ as well
as discrete indices. Repeated labels are summed and integrated. All quantum
variables are understood to be evaluated at the same time, with the common
time argument dropped everywhere. The quantities $\mathcal{Q}^r$ are the
same as the $C^r$ introduced in Section 7.2.} one set $Q^n$ of independent
canonical variables (like $V^i(\mathbf{x})$) with independent canonical
conjugates $P_n=\partial L/\partial\dot Q^n$; and another set
$\mathcal{Q}^r(\mathbf{x})$ (like $V^0$) whose time-derivatives do not appear
in the Lagrangian. The primary constraints are the conditions
$\chi_{1r}=0$, where
\begin{equation}
  \chi_{1r}=\mathcal{P}_r
  \tag{7.A.1}\label{eq:7.A.1}
\end{equation}
are the variables conjugate to the $\mathcal{Q}^r$. The secondary constraints
arise from the equations of motion $0=\partial L/\partial\mathcal{Q}^r$ for
the $\mathcal{Q}^r$; we suppose that these constraints can be `solved'---i.e.,
they may be written in the form $\chi_{2r}=0$, with $\chi_{2r}$ in the form
\begin{equation}
  \chi_{2r}=\mathcal{Q}^r-f^r(Q,P).
  \tag{7.A.2}\label{eq:7.A.2}
\end{equation}
(An example is provided by Eq.~\eqref{eq:7.6.5}, which gives $V^0$ in terms of the
independent $P$s (here, the $\Pi_i$) and $Q$s.) We assume that the independent
$Q$s and $P$s satisfy the usual canonical commutation rules:
\begin{equation}
  [Q^n,P_m]=i\delta^n_m,\qquad
  [Q^n,Q^m]=[P_n,P_m]=0.
  \tag{7.A.3}\label{eq:7.A.3}
\end{equation}
The constraint $\chi_{2r}=0$ yields the commutators involving
$\mathcal{Q}^r$:
\begin{equation}
  [\mathcal{Q}^r,Q^n]=-i\frac{\partial f^r}{\partial P_n},\qquad
  [\mathcal{Q}^r,P_n]=i\frac{\partial f^r}{\partial Q^n},
  \tag{7.A.4}\label{eq:7.A.4}
\end{equation}
\begin{equation}
  [\mathcal{Q}^r,\mathcal{Q}^s]=i\Gamma^{rs},
  \tag{7.A.5}\label{eq:7.A.5}
\end{equation}
where $\Gamma^{rs}$ is the Poisson bracket
\begin{equation}
  \Gamma^{rs}\equiv[f^r,f^s]_{\mathrm P},
  \tag{7.A.6}\label{eq:7.A.6}
\end{equation}
and, of course, all commutators involving $\mathcal{P}_r$ vanish:
\begin{equation}
  [\mathcal{P}_r,Q^n]=[\mathcal{P}_r,P_n]
  =[\mathcal{P}_r,\mathcal{Q}^s]=[\mathcal{P}_r,\mathcal{P}_s]=0.
  \tag{7.A.7}\label{eq:7.A.7}
\end{equation}

Now let us compare these commutators with the Dirac brackets. The Poisson
brackets of the constraint functions are
\begin{equation}
  C_{1r,1s}\equiv[\chi_{1r},\chi_{1s}]_{\mathrm P}=0,
  \tag{7.A.8}\label{eq:7.A.8}
\end{equation}
\begin{equation}
  C_{1r,2s}=-C_{2s,1r}\equiv[\chi_{1r},\chi_{2s}]_{\mathrm P}
  =-\delta^s_r,
  \tag{7.A.9}\label{eq:7.A.9}
\end{equation}
\begin{equation}
  C_{2r,2s}\equiv[\chi_{2r},\chi_{2s}]_{\mathrm P}
  =[f^r(Q,P),f^s(Q,P)]_{\mathrm P}=\Gamma^{rs}.
  \tag{7.A.10}\label{eq:7.A.10}
\end{equation}
(In the example of the massive vector field $\Gamma^{rs}$ vanishes, but the
discussion here will apply also for non-vanishing $\Gamma^{rs}$.) It is easy
to see that the $C$-matrix has the inverse
\begin{equation}
  \begin{gathered}
    (C^{-1})^{1r,1s}=\Gamma^{rs},\qquad
    (C^{-1})^{2r,2s}=0,\\
    (C^{-1})^{1r,2s}=-(C^{-1})^{2s,1r}=\delta^r_s.
  \end{gathered}
  \tag{7.A.11}\label{eq:7.A.11}
\end{equation}
Also, the Poisson brackets of any function $A$ with the constraint functions
are
\[
  [A,\chi_{1r}]_{\mathrm P}=\frac{\partial A}{\partial\mathcal{Q}^r},
  \qquad
  [A,\chi_{2r}]_{\mathrm P}
  =-\frac{\partial A}{\partial\mathcal{P}_r}-[A,f^r(Q,P)]_{\mathrm P}.
\]
Hence the Dirac bracket is
\begin{equation}
  \begin{aligned}
    [A,B]_{\mathrm D}={}&[A,B]_{\mathrm P}
    -\frac{\partial A}{\partial\mathcal{Q}^r}
      \frac{\partial B}{\partial\mathcal{P}_r}
    +\frac{\partial B}{\partial\mathcal{Q}^r}
      \frac{\partial A}{\partial\mathcal{P}_r}\\
    &+\frac{\partial A}{\partial\mathcal{P}_r}\Gamma^{rs}
      \frac{\partial B}{\partial\mathcal{P}_s}
    -\frac{\partial A}{\partial\mathcal{P}_r}[B,f^r]_{\mathrm P}
    +[A,f^r]_{\mathrm P}\frac{\partial B}{\partial\mathcal{P}_r}.
  \end{aligned}
  \tag{7.A.12}\label{eq:7.A.12}
\end{equation}
Now, if $A$ and $B$ are both functions only of the independent canonical
variables $Q^n$ and $P_n$, then
$\partial A/\partial\mathcal{Q}^r=\partial B/\partial\mathcal{Q}^r=0$, so the
Dirac bracket is equal to the Poisson bracket. In particular,
\begin{equation}
  [Q^n,P_m]_{\mathrm D}=\delta^n_m,\qquad
  [Q^n,Q^m]_{\mathrm D}=[P_n,P_m]_{\mathrm D}=0.
  \tag{7.A.13}\label{eq:7.A.13}
\end{equation}
Where $A$ is $\mathcal{Q}^r$ and $B$ is a function of $Q$s and $P$s, it is
only the fifth term on the right-hand side of Eq.~\eqref{eq:7.A.12} that contributes.
In particular
\begin{equation}
  [\mathcal{Q}^r,Q^n]_{\mathrm D}=-\frac{\partial f^r}{\partial P_n},
  \qquad
  [\mathcal{Q}^r,P_n]_{\mathrm D}=+\frac{\partial f^r}{\partial Q^n}.
  \tag{7.A.14}\label{eq:7.A.14}
\end{equation}
Where both $A$ and $B$ are $\mathcal{Q}$s, we have only the fourth term
\begin{equation}
  [\mathcal{Q}^r,\mathcal{Q}^s]_{\mathrm D}=\Gamma^{rs}.
  \tag{7.A.15}\label{eq:7.A.15}
\end{equation}
Finally, where $A$ is $\mathcal{P}_r$ and $B$ is anything, we have only the
first and third terms, which cancel:
\begin{equation}
  [\mathcal{P}_r,B]_{\mathrm D}
  =[\mathcal{P}_r,B]_{\mathrm P}
  +\frac{\partial B}{\partial\mathcal{Q}^r}=0.
  \tag{7.A.16}\label{eq:7.A.16}
\end{equation}
Comparison of Eqs.~\eqref{eq:7.A.13}--\eqref{eq:7.A.16} with Eqs.~\eqref{eq:7.A.3}--\eqref{eq:7.A.7} shows that
in all cases, the commutators are equal to the Dirac brackets times $i$. This
is only to be expected, because as remarked in Section 7.6 all Dirac brackets
involving the constraint functions vanish, so the Dirac brackets involving
$\mathcal{Q}^r$ and/or $\mathcal{P}_s$ are given by using the constraint
equations to express $\mathcal{Q}^r$ and/or $\mathcal{P}_s$ in terms of the
independent $Q$s and $P$s.

\begin{center}
$B$
\end{center}

Next consider the case where the constraints take the form of conditions
$\chi_{1r}(\Psi)=0$ on the $\Psi^a$, which can be solved by expressing them in
terms of a smaller set of unconstrained variables $Q^n$, and an equal number
of separate conditions $\chi_{2r}(\Pi)=0$ on the $\Pi_a$, which can be solved
by expressing the $\Pi_a$ in terms of a smaller set of unconstrained $P_n$.
(We will see an example in the next chapter, where the constraints on the
$\Psi^a$ are gauge fixing conditions that are used to eliminate first class
constraints, and the constraints on the $\Pi_a$ are secondary constraints
arising from the consistency of the first class constraints with the field
equations.) We assume that the unconstrained variables satisfy the usual
canonical commutation relations
$[Q^n,P_m]=i\delta^n_m$, $[Q^n,Q^m]=[P_n,P_m]=0$. The constrained and
unconstrained momenta are related by
\begin{equation}
  P_n=\frac{\partial L}{\partial\dot Q^n}
  =\frac{\partial L}{\partial\dot\Psi^b}
    \frac{\partial\dot\Psi^b}{\partial\dot Q^n}
  =\Pi_b\frac{\partial\Psi^b}{\partial Q^n}.
  \tag{7.A.17}\label{eq:7.A.17}
\end{equation}
It follows that
\[
  [\Psi^a,\Pi_b]\frac{\partial\Psi^b}{\partial Q^n}
  =[\Psi^a,P_n]=i\frac{\partial\Psi^a}{\partial Q^n}
\]
or, in other words,
\begin{equation}
  \left\{[\Psi^a,\Pi_b]-i\delta^a_b\right\}
  \frac{\partial\Psi^b}{\partial Q^n}=0.
  \tag{7.A.18}\label{eq:7.A.18}
\end{equation}
Now, the constraint $\chi_{1r}(\Psi)=0$ is satisfied for
$\Psi^a=\Psi^a(Q)$ for all $Q$, so
\begin{equation}
  \frac{\partial\chi_{1r}}{\partial\Psi^b}
  \frac{\partial\Psi^b}{\partial Q^n}=0.
  \tag{7.A.19}\label{eq:7.A.19}
\end{equation}
Furthermore, the vectors $(V_r)_b\equiv\partial\chi_{1r}/\partial\Psi^b$
form a complete set perpendicular to all the vectors
$(U_n)^b\equiv\partial\Psi^b/\partial Q^n$, because if there were some other
vector $V_b$ with $V_b(U_n)^b=0$ for all $n$, then there would be additional
constraints on the $\Psi^a$. Hence Eq.~\eqref{eq:7.A.18} implies that
\begin{equation}
  [\Psi^a,\Pi_b]=i\delta^a_b+ic^a_r
  \frac{\partial\chi_{1r}}{\partial\Psi^b}
  \tag{7.A.20}\label{eq:7.A.20}
\end{equation}
with some unknown coefficients $c^a_r$. To determine these coefficients, we
make use of the other constraint, that $\chi_{2r}(\Pi)=0$. It follows that
\[
  0=[\Psi^a,\chi_{2r}(\Pi)]
  =[\Psi^a,\Pi_b]\frac{\partial\chi_{2r}(\Pi)}{\partial\Pi_b}.
\]
Using Eq.~\eqref{eq:7.A.20}, we have then
\begin{equation}
  \frac{\partial\chi_{2r}(\Pi)}{\partial\Pi_a}
  =-c^a_s\frac{\partial\chi_{1s}(\Psi)}{\partial\Psi^b}
    \frac{\partial\chi_{2r}(\Pi)}{\partial\Pi_b}.
  \tag{7.A.21}\label{eq:7.A.21}
\end{equation}
We recognize the factor multiplying $c^a_s$ as the Poisson bracket
\[
  \frac{\partial\chi_{1s}(\Psi)}{\partial\Psi^b}
  \frac{\partial\chi_{2r}(\Pi)}{\partial\Pi_b}
  =[\chi_{1s},\chi_{2r}]_{\mathrm P}\equiv C_{1s,2r}.
\]
Also, since $\chi_{1s}$ depends only on the $\Psi$ and $\chi_{2r}$ depends
only on the $\Pi$, these are the only non-vanishing Poisson brackets of
constraints, so
\[
  C_{1r,1s}=C_{2r,2s}=0.
\]
Thus Eq.~\eqref{eq:7.A.21} may be written
\begin{equation}
  \frac{\partial\chi_N}{\partial\Pi_a}=-c^a_s C_{1s,N}
  \tag{7.A.22}\label{eq:7.A.22}
\end{equation}
with $N$ running over all constraint functions. For second class constraints,
this has the unique solution
\begin{equation}
  c^a_s=-\frac{\partial\chi_N}{\partial\Pi_a}(C^{-1})^{N,1s}
  =-\frac{\partial\chi_{2r}}{\partial\Pi_a}(C^{-1})^{2r,1s}.
  \tag{7.A.23}\label{eq:7.A.23}
\end{equation}
Using this in Eq.~\eqref{eq:7.A.20} shows that
\begin{equation}
  [\Psi^a,\Pi_b]
  =i\left[\delta^a_b
  -\frac{\partial\chi_{2r}}{\partial\Pi_a}(C^{-1})^{2r,1s}
   \frac{\partial\chi_{1s}}{\partial\Psi^b}\right].
  \tag{7.A.24}\label{eq:7.A.24}
\end{equation}
The Poisson brackets of $\Psi^a$ and $\Pi_b$ with the constraint functions are
\begin{equation}
  \begin{gathered}
    [\Psi^a,\chi_{1r}]_{\mathrm P}=0,
    \qquad
    [\Psi^a,\chi_{2r}]_{\mathrm P}
    =\frac{\partial\chi_{2r}}{\partial\Pi_a},\\
    [\chi_{1r},\Pi_b]_{\mathrm P}
    =\frac{\partial\chi_{1r}}{\partial\Psi^b},
    \qquad
    [\chi_{2r},\Pi_b]_{\mathrm P}=0,
  \end{gathered}
  \tag{7.A.25}\label{eq:7.A.25}
\end{equation}
so the quantity in brackets on the right-hand side of Eq.~\eqref{eq:7.A.24} is the
Dirac bracket
\begin{equation}
  [\Psi^a,\Pi_b]=i[\Psi^a,\Pi_b]_{\mathrm D}
  \tag{7.A.26}\label{eq:7.A.26}
\end{equation}
as was to be shown. Also, we can easily see that because $C^{-1}$ has no 11
or 22 components, the other Dirac brackets are
\begin{equation}
  [\Psi^a,\Psi^b]_{\mathrm D}=[\Pi_a,\Pi_b]_{\mathrm D}=0,
  \tag{7.A.27}\label{eq:7.A.27}
\end{equation}
so trivially
\begin{equation}
  [\Psi^a,\Psi^b]=i[\Psi^a,\Psi^b]_{\mathrm D},
  \qquad
  [\Pi_a,\Pi_b]=i[\Pi_a,\Pi_b]_{\mathrm D}.
  \tag{7.A.28}\label{eq:7.A.28}
\end{equation}

\begin{center}
$*\quad*\quad*$
\end{center}

In addition to commutation rules, we also need an explicit formula for the
Hamiltonian. The usual canonical formalism tells us to take
\begin{equation}
  H=P_n\dot Q^n-L,
  \tag{7.A.29}\label{eq:7.A.29}
\end{equation}
the sum running over independent canonical variables. In theories of both
types considered in this Appendix, this Hamiltonian may be written in terms
of the constrained variables as
\begin{equation}
  H=\Pi_a\dot\Psi^a-L.
  \tag{7.A.30}\label{eq:7.A.30}
\end{equation}
For theories of type A, this is trivial; the sum over $a$ runs over values
$n$, for which $\Psi^n=Q^n$ and $\Pi_n=P_n$ are the independent canonical
variables, together with values $r$, for which
$\Pi_r=\mathcal{P}_r=0$. For theories of type B, we note that Eq.~\eqref{eq:7.A.17}
gives
\[
  P_n\dot Q^n=\Pi_b\frac{\partial\Psi^b}{\partial Q^n}\dot Q^n
  =\Pi_b\dot\Psi^b
\]
which again yields Eq.~\eqref{eq:7.A.30}.
